% =========================================================
% Proof: Subsequences of Convergent Sequences Converge to the Same Limit
% =========================================================

\newpage
\phantomsection
\label{prf:subsequence-convergence-q}

\begin{remark*}[Return]
\hyperref[thm:subsequence-convergence-q]{Return to Theorem~\ref*{thm:subsequence-convergence-q}.}
\end{remark*}

\begin{theorem*}[Subsequences of convergent sequences converge to the same limit]
Let $(x_n)$ be a sequence of rational numbers with $x_n \to a \in \mathbb{Q}$.
Then every subsequence $(x_{k_n})$ converges to $a$.
\end{theorem*}

\begin{proof}[Proof (Professional Standard)]
Let $(x_{k_n})$ be any subsequence of $(x_n)$.  Let $\varepsilon > 0$.
Since $x_n \to a$, there exists $N \in \mathbb{N}$ such that
$|x_n - a| < \varepsilon$ for all $n \geq N$.
We claim $k_n \geq n$ for all $n \in \mathbb{N}$.
Indeed, $k_1 \geq 1$ since $(k_n) \subseteq \mathbb{N}$, and if
$k_n \geq n$ then $k_{n+1} > k_n \geq n$, so $k_{n+1} \geq n+1$.
By induction, $k_n \geq n$ for all $n$.
Therefore, for all $n \geq N$ we have $k_n \geq n \geq N$, hence
$|x_{k_n} - a| < \varepsilon$.
Since $\varepsilon > 0$ was arbitrary, $x_{k_n} \to a$.
\end{proof}

\begin{proof}[Proof (Detailed Learning)]
\textbf{Step 1: Establish the key lemma $k_n \geq n$.}
We prove by induction that $k_n \geq n$ for all $n \in \mathbb{N}$.

\emph{Base case:} $k_1 \in \mathbb{N}$, so $k_1 \geq 1$.

\emph{Inductive step:} Suppose $k_n \geq n$ for some $n$.
Since the sequence $(k_n)$ is strictly increasing, $k_{n+1} > k_n$.
Since $k_{n+1}$ is a natural number and $k_{n+1} > k_n \geq n$,
we get $k_{n+1} \geq n + 1$.

By induction, $k_n \geq n$ for all $n \in \mathbb{N}$.

\textbf{Step 2: Use the convergence of $(x_n)$.}
Let $\varepsilon > 0$.  Since $x_n \to a$, there exists $N \in \mathbb{N}$
such that $|x_n - a| < \varepsilon$ for all $n \geq N$.

\textbf{Step 3: Verify the subsequence condition.}
For any $n \geq N$: by Step~1, $k_n \geq n \geq N$.
Therefore $|x_{k_n} - a| < \varepsilon$ by the choice of $N$.

\textbf{Step 4: Conclude.}
Since $\varepsilon > 0$ was arbitrary, we have shown that for each
$\varepsilon > 0$ there is an $N$ such that $|x_{k_n} - a| < \varepsilon$
for all $n \geq N$.  By definition, $x_{k_n} \to a$.
\end{proof}

\begin{remark*}[Proof structure]
The structural key is the lemma $k_n \geq n$: it ensures that every
subsequence index eventually exceeds any given threshold $N$, allowing
us to transfer the convergence bound from $(x_n)$ to $(x_{k_n})$.
The induction is the only non-trivial step; everything else is definition-unwinding.
\end{remark*}

\begin{remark*}[Dependencies]
\hyperref[def:convergent-rational-sequence]{Convergence of a rational sequence},
\hyperref[def:subsequence-q]{Subsequence},
induction on $\mathbb{N}$.
\end{remark*}
